\appendix
\chapter{Implementation}
\section{se\_vault configuration example}

\begin{lstlisting}[breaklines=true]
<start name="se_vault" caps="1500" ram="2000M">
 <exit propagate="yes"/>
 <provides>
  <service name="File_system"/>
 </provides>
 <config
 jitterentropy_available="yes"
 fsck_apply="yes">
  <vfs/>
 </config>
 <route>
  <service name="Block"
  label_prefix="se_tresor_trust_anchor_vfs ->"> 
   <child name="usb_sec"
   identity="se_vault -> trust_anchor"/> 
  </service>
  <service name="Block"
  label_prefix="se_tresor_init ->">
   <child name="nvme2"
   identity="se_vault -> data"/>
  </service>
  <service name="Block"
  label_prefix="tresor ->">
   <child name="nvme2"
   identity="se_vault -> data"/>
  </service>
  <service name="Block"
  label_prefix="se_tresor_vfs ->">
   <child name="nvme2"
   identity="se_vault -> data"/>
  </service>
  <service name="Block"
  label_prefix="init_flag ->">
   <child name="nvme2"
   identity="se_vault -> data"/>
  </service>
  <service name="Timer">
   <child name="timer"/>
  </service>
  <service name="Report">
   <child name="report_rom"/>
   </service>
  <service name="PD">
   <parent/>
  </service>
  <service name="ROM">
   <parent/>
  </service>
  <service name="CPU">
   <parent/>
  </service>
  <service name="LOG">
   <parent/>
  </service>
  <service name="RM">
   <parent/>
  </service>
  <any-service>
   <parent/>
  </any-service>
 </route>
</start>
\end{lstlisting}
